%@ pages 209-222
%@ status translated
% Original section 19, "La descente pour les préfaisceaux de Quillen à gauche", pp. 209--222.
%
% Translation notes (the mathematics is left exactly as printed):
%  * "cmf" is rendered "cmc", as throughout this translation; "préchamp"
%    and "préchamp de Segal" are "prestack" and "Segal prestack"; "champ"
%    is "stack" and "protochamp" is "protostack" (see \S10).
%  * "section eq" is kept as "eq-section", the authors' abbreviation.
%  * Four slips of the original are kept as printed and marked with a
%    % comment at the spot: the reference to $(o)''$ where the type (II)
%    structure comes from $(o)'''$; the right adjoint called
%    $r_{\Uu}^{\ast}$ instead of $r_{\Uu,\ast}$ in the proof of
%    Theorem \ref{dansleschamps}; $M(X)$ for $\mm(X)$ in the same
%    proof; and the composite $\epsilon^{\ast}C_{a,-}\to
%    \epsilon^{\ast}C_{a,-}$ of Illusie 1.1.5, whose source should
%    presumably be $\epsilon^{\ast}\sigma(a)$. None of them affects the
%    mathematics.
%  * The v1/v2/v3 comments of the source are kept, translated.

\clearpage
\section{Descent for left Quillen presheaves}
\label{quidescentepage}

We shall prove a descent theorem for a left Quillen presheaf, which will
give conditions for $L(\mm )$ to be a stack. We found this proof starting
from the case of complexes. The main problem for this result is that one
must start from a weak section of $L(\mm )$ which may, {\em a priori}, be
rather ``wild''. In the preceding chapters we have found a series of results
which allow one to ``tame'' such a section, i.e.\ to reduce to the case of a
section of $\int \mm$ instead of a section of $\int 'L(\mm )$.

The work now consists in proving that a section of $\int \mm$ descends to
the base object. This last part is in reality already well known; for
example, one knows how to take the simple complex associated to a
cosimplicial complex, see Deligne \cite{HodgeIII}. We conjecture that it is
this problem which was addressed in the ``secret papers'' of Deligne to
which Illusie refers in \cite{Illusie}. It is also the subject of the Exposé
of B.~Saint-Donat on {\em cohomological descent} in SGA 4 (\cite{SGA4}). The
method in all cases is to start with a section $\sigma$ in
$Sect ^{\rm eq} (\Delta , \int _{\Delta} \mm )$ (in SGA 4 this set of
sections is called $\Gamma ^{\rm cocart}(\, \, )$) and to ``descend'' it via
the augmentation, by applying the ``direct image'' functor which we shall
denote $\varphi _{\ast}$. The problem is to verify that this really answers
the descent problem, i.e.\ that
$\varphi ^{\ast} (\varphi _{\ast} \sigma )$ is equivalent to $\sigma$.

See the end of this chapter for a more detailed comparison between our
discussion and that of SGA 4.

At the beginning of this section we shall study a left Quillen presheaf
$\mm$ on $(\Delta ^+)^o$. If $p\rightarrow q$ is a morphism of $\Delta^+$
then the restriction functor (i.e.\ the left Quillen functor) goes in the
direction $\mm (p)\rightarrow \mm (q)$. Here $\Delta ^+$ is the category
$\Delta$ augmented by an initial object which we shall denote $\iota$.
Except in case of possible confusion, we shall avoid the notation
$\mm |_{\Delta ^o}$ and shall also denote this restriction by $\mm$.

We shall make the following hypothesis (o):
\newline
$(o)'$\,\, that each $\mm (y)$ admits arbitrary small limits and colimits,
and functorial factorisations;
\newline
$(o)''$\,\,
that every object of $\mm (y)$ is cofibrant; and
\newline
$(o)'''$\,\, that $Sect (\Delta ^+, \int
_{\Delta ^+}\mm )$ and $Sect (\Delta , \int _{\Delta} \mm )$ admit
cmc structures of type (II)
(of Theorem \ref{sectionscmf}) generated by
cofibrations.

Hypothesis $(o)''$, that every object is cofibrant, is there purely for
convenience and is obviously inessential. The structure of type (II) of
hypothesis $(o)'''$ could probably be replaced, in our argument, by a Reedy
structure. On the other hand, hypothesis $(o)'$ (which moreover has become
standard, cf.\ \cite{Hirschhorn} \cite{DHK} \cite{HoveyBook}) and in
particular the closure of $\mm (\iota )$ under small limits---for example,
under limits indexed by $\Delta$---is essential, and the theorem
\ref{dansleschamps} below would probably no longer be true without this
condition.

% the original refers to $(o)''$ here; the type (II) structure in question
% is the one of $(o)'''$.
We equip $Sect (\Delta , \int _{\Delta} \mm )$ with its
cmc structure of type (II) from hypothesis $(o)''$.
The ``constant section'' functor
$$
\varphi ^{\ast} : \mm (\iota )\rightarrow Sect (\Delta , \int _{\Delta} \mm )
$$
is then a left Quillen functor, and its right adjoint, which we shall denote
$\varphi _{\ast}$ (which exists by hypothesis $(o)'$), is a right Quillen
functor.

One can decompose $\varphi _{\ast}$ in the following way. The adjoint of the
restrictions provides a morphism
%%%
%%% changed from r to r_{\ast} in v2
%%%
$$
r_{\ast}:Sect (\Delta , \int _{\Delta} \mm )\rightarrow \mm (\iota )^{\Delta}.
$$
Composing
next with the functor $lim : \mm (\iota )^{\Delta} \rightarrow \mm (\iota
)$ one obtains
$$
\varphi _{\ast} : Sect (\Delta , \int _{\Delta} \mm )\rightarrow \mm (\iota ),
$$
$$
\varphi _{\ast} (\sigma ) := \lim _{\Delta} r_{\ast}\sigma .
$$

Note that a fibrant object for the structure of type (II) is also fibrant
for the Reedy structure (which exists automatically on sections).
Moreover $r$ takes fibrant objects to fibrant objects
(it is a right Quillen functor, being a composite of adjoints of
restrictions). Hence, for $\sigma \in Sect (\Delta , \int _{\Delta} \mm )_f$,
$\varphi _{\ast} (\sigma )$ is equivalent to $holim _{\Delta}r_{\ast}\sigma $.


We want to show that $\varphi ^{\ast}$ induces an equivalence
$$
L(\varphi ^{\ast} ): L(\mm (\iota ))\cong
L(Sect ^{\rm eq}(\Delta , \int _{\Delta} \mm )).
$$
We begin with the following observation. For this lemma we keep the
convenience hypothesis that all objects are cofibrant, but it is certainly
not essential.


\begin{lemma}
\label{dpq1}
Let $N$ and $M$ be closed model categories and
$\varphi _{\ast} : N\rightarrow M$, $\varphi ^{\ast} : M\rightarrow N$ a
pair of adjoint functors,
where $\varphi ^{\ast}$ is a left Quillen functor and
$\varphi _{\ast}$ a
right Quillen functor. Suppose that all objects of $M$ and
of $N$ are cofibrant. Suppose that for every object $x\in M$, the adjunction
morphism $x\rightarrow \varphi _{\ast} (\varphi ^{\ast} x)'$ is an
equivalence (where $(\varphi ^{\ast}
x)'$ is the fibrant replacement of $\varphi ^{\ast} x$).
\newline Then the functor
$$
L(\varphi ^{\ast} ) : L(M)\rightarrow L(N)
$$
is fully faithful, with essential image
consisting of the fibrant objects $y\in N$
such that $\varphi ^{\ast} (\varphi _{\ast} y)\rightarrow y$ is an
equivalence, and with
$L(\varphi _{\ast} )$ as inverse on this image.
\end{lemma}
{\em Proof:}
Let $N'\subset N_f$ be the subcategory of the fibrant objects $y$ such that
$\varphi ^{\ast} (\varphi _{\ast} y)\rightarrow y$
is an equivalence. Note that $N'$ is stable under equivalence in $N_f$
(i.e.\ if $z\in N_f$
and $y\in N'$ with $z\cong y$ then $z\in N'$). By the principle \ref{stabilite},
$(N',W_{N_f}\cap N')$ admits a homotopical calculus of fractions, and the
localised category
$L(N',W_{N_f}\cap N')$ is a full simplicial subcategory of
$L(N_f,W_{N_f})\cong L(N,W_N)$. By \cite{DwyerKan2} Corollary 3.6, the
functors $\varphi _{\ast}$ and $\varphi ^{\ast}$ establish an
equivalence between $L(M, W_M)$
and $L(N', W_N\cap N')$.
\eop


Let us return to the preceding situation. We shall apply the lemma to
the pair of
functors $\varphi _{\ast} , \varphi ^{\ast}$ defined previously.
If $\sigma \in
Sect (\Delta , \int _{\Delta} \mm )$ is equivalent to $\varphi ^{\ast}
(x)$ for $x\in
\mm (\iota )$ then the morphisms $\sigma ^r(f)$ are equivalences, i.e.\
$\sigma$ belongs to $Sect ^{\rm eq} (\Delta , \int _{\Delta} \mm )$.

We shall keep hypothesis (o) above.
In addition, we shall suppose that we have the two following
properties: \newline
(i)\,\, a morphism $a$ of $\mm (\iota )$ is an equivalence
if and only if $\varphi ^{\ast} (a)$ is an equivalence; and
\newline
(ii)\,\, for every $\sigma \in Sect ^{\rm eq} (\Delta , \int _{\Delta} \mm )_f$,
the adjunction morphism $\varphi ^{\ast} \varphi _{\ast} (\sigma )
\rightarrow \sigma$
is an equivalence.

One may remark that property (i) is equivalent
to the condition:
\newline
(i)'\,\, a morphism $a$ of $\mm (\iota )$ is
an equivalence if and only if
its restriction to $\mm (0)$ is one.

Then for $x\in \mm (\iota )$ property (ii) applies to
$\varphi ^{\ast} (x)'$
(the fibrant replacement of $\varphi ^{\ast} (x)$),
so the morphism
$\varphi ^{\ast} \varphi _{\ast} (\varphi ^{\ast} (x)')\rightarrow \varphi
^{\ast} (x)'$ is an equivalence. The
composite
$$
\varphi ^{\ast} (x)\rightarrow \varphi ^{\ast}
 \varphi _{\ast} (\varphi ^{\ast} (x)') \rightarrow \varphi ^{\ast} (x)'
$$
is equal to the weak equivalence $\varphi ^{\ast} (x)\rightarrow
\varphi ^{\ast} (x)'$ of the
fibrant replacement. The first morphism is therefore a weak equivalence,
but this morphism is $\varphi ^{\ast} (a)$ where $a$ is the adjunction
morphism
$a:x\rightarrow \varphi _{\ast} \varphi ^{\ast} (x)$. The hypothesis of the
lemma is therefore verified; it follows that
$L(\varphi ^{\ast} )$ is fully faithful. On the other hand,
condition (ii)
above
identifies the image of $L(\varphi ^{\ast} )$ as the localised category of
$Sect
^{\rm eq} (\Delta
, \int _{\Delta} \mm )$.

We note as above that (i) is equivalent to the condition
according to which
a morphism $a$ of $\mm (\iota )$ is an equivalence
if and only if its restriction to $\mm (0)$ is one.
We have obtained the

\begin{corollary}
\label{dpq2}
Let $\mm$ be a left Quillen presheaf on $(\Delta ^+)^o$
which satisfies the conditions above, which we recall:
\newline
(o)\,\, each $\mm (y)$ admits arbitrary small limits and colimits,
and functorial factorisations;
every object of $\mm (y)$ is cofibrant; and $Sect (\Delta ^+, \int
_{\Delta ^+}\mm )$ and $Sect (\Delta , \int _{\Delta} \mm )$ admit
cmc structures of type (II) (of
Theorem \ref{sectionscmf}) generated by
cofibrations.
\newline
(i)\,\, a morphism $a$ of $\mm (\iota )$ is an equivalence
if and only if its restriction to $\mm (0)$ is one; and
\newline
(ii)\,\, for every $\sigma \in Sect ^{\rm eq} (\Delta , \int _{\Delta} \mm )_f$,
the adjunction morphism $\varphi ^{\ast} \varphi _{\ast} (\sigma )
\rightarrow \sigma$
is an equivalence.
\newline
Then the morphism
$$
L(\mm (\iota ))\rightarrow
L(Sect ^{\rm eq} (\Delta , \int _{\Delta} \mm ))
$$
is an equivalence of simplicial categories.
\end{corollary}
\eop

\bigskip

We now take up the case of a Quillen presheaf
on a site. Let $\Xx$ be a site admitting fibre products and sufficiently
many disjoint sums (see \S 15) and let $\mm$ be a left Quillen presheaf
on $\Xx$.
If $f:X\rightarrow Y$ is a
morphism of $\Xx$ then we denote by
$$
f^{\ast}: \mm (Y)\rightarrow \mm (X)
$$
the restriction $res_f$, and by $f_{\ast}$ its right adjoint.

We make the hypothesis that $L(\mm )$ is compatible with disjoint sums
\ref{compsomdis}. We first present a criterion which allows one to verify
this hypothesis.
If $\Uu = \{ U_{\alpha}\}$ is a family of objects of $\Xx$, we set
$$
\mm (\Uu ):= \prod _{\alpha} \mm (U_{\alpha}).
$$
This has a cmc structure.
The restriction morphism provides a morphism
$$
r_{\Uu}^{\ast}:\mm (\amalg \Uu )\rightarrow \mm (\Uu ),
$$
which is a
left Quillen functor,
whose right adjoint we denote by $r_{\Uu , \ast}$.

\begin{corollary}
\label{compsomcritere}
Suppose that, for every family $\Uu$ of objects of $\Xx$ (of size $<\beta$),
and for every $u$ in $\mm (\Uu )$, the adjunction morphism
$$
r_{\Uu}^{\ast} r_{\Uu , \ast} (u)\rightarrow u
$$
is a weak equivalence. Suppose moreover that, for every family $\Uu$
as above, a
morphism $a$ of $\mm (\amalg \Uu )$ is a weak equivalence
if and only if $r_{\Uu}^{\ast} (a)$ is one. Then
$L(\mm )$ is compatible with disjoint sums.
\end{corollary}
{\em Proof:} This is a well-known characterisation of {\em Quillen
equivalences}. Moreover, it is a corollary of Lemma \ref{dpq1}.
\eop

%%% sentence added in v2/v3
The statement of the following theorem, already corrected in
version 2 of the paper by the addition of hypothesis (4), is corrected in the
present version 3 by the addition of
hypothesis (5).


\begin{theorem}
\label{dansleschamps}
Suppose
that $\Xx$ is a site which admits fibre products and sufficiently many
disjoint sums compatible with fibre products.
Let
$\mm$ be a Quillen presheaf on $\Xx$. We suppose that $\mm$ satisfies the
three following properties: \newline
(0)\,\, each $\mm (y)$ admits arbitrary small limits and colimits,
and functorial factorisations;
every object of $\mm (y)$ is cofibrant; and
for every functor $\Yy \rightarrow \Xx$ of $1$-categories,
$Sect (\Yy , \int _{\Yy }\mm |_{\Yy})$ admits a cmc structure of
type (II) generated by cofibrations.
\newline
(1)\,\, For every $X$ of $\Xx$ and
every family
$\Uu =\{ U_{\alpha}\stackrel{p_{\alpha}}{\rightarrow}
X\}$ covering $X$ (or satisfying $X=\amalg \Uu$)
a morphism $a$ of $\mm (X)$ is a weak equivalence
if and only if its restriction $p_{\alpha}^{\ast}(a)$ to $\mm
(U_{\alpha})$ is a weak equivalence for every $\alpha$;
\newline
(2)\,\, if $X= \amalg \Uu$ is the disjoint sum of a family $\Uu =\{
U_{\alpha}\}$ then the adjunction morphism
$$
r_{\Uu}^{\ast} r_{\Uu , \ast}u\rightarrow u
$$
is a weak equivalence for every $u\in \mm (\Uu ):= \prod _{\alpha} \mm
(U_{\alpha})$;
\newline
(3)\,\, the presheaf $\mm$ is {\em cartesian} in the sense that for
every cartesian diagram of objects of $\Xx$
$$
\begin{array}{ccc}
Y\times _XZ & \stackrel{q}{\rightarrow} & Z \\
{\scriptstyle p}\downarrow && \downarrow {\scriptstyle r}\\
Y& \stackrel{s}{\rightarrow} & X
\end{array}
$$
and for every $u$ in $\mm (Y)$.
the natural map
$$
r^{\ast} s_{\ast} (u) \stackrel{\cong}{\rightarrow} q_{\ast} p^{\ast}(u).
$$
is a weak
equivalence;
%%% condition added in v2
\newline
(4)\,\, the restriction morphisms of the localised prestack $L(\mm )$
preserve homotopy limits (note that these morphisms preserve
homotopy colimits automatically, since
$\mm$ is left Quillen); and
%%%%%%%% condition (5) added in v3 : %%%%%%%%%%%
\newline
(5)\,\, the Segal prestack $L(\mm )$ is a protostack, that is, for every
$x\in \Xx$ and all $m,m' \in L(\mm )(x)$ the simplicial
presheaf $L(\mm )_{1/}(m,m')$ is a stack on $\Xx / x$.


Then the presheaf of simplicial categories $L(\mm )$ is a stack (i.e.\ a
Segal $1$-stack) on $\Xx$.
\end{theorem}
{\em Proof:}
We first show, with the help of Corollary \ref{compsomcritere},
that $L(\mm )$ is compatible with disjoint sums. If $X=\amalg U_{\alpha}$
is the disjoint sum of the family $\Uu = \{ U_{\alpha}\}$ then
we have the functor
$$
r_{\Uu}^{\ast}: \mm (X)\rightarrow \prod _{\alpha} \mm (U_{\alpha})
$$
% the original prints $r_{\Uu}^{\ast}$ for the right adjoint here
and its right adjoint $r_{\Uu}^{\ast}$. Condition (1) together with the fact
that $\Uu$ is a family covering $X$ imply that an arrow $a$ of
$\mm (X)$ is a weak equivalence if and only if $r^{\ast}_{\Uu}(a)$
is a weak equivalence. Condition (2) gives the other part of
the hypothesis of Corollary \ref{compsomcritere},
and we therefore obtain that $L(\mm )$
is compatible with disjoint sums.

We shall apply \ref{critere7} to $L(\mm )$. We may therefore fix $X\in \Xx$
and a covering family $\Uu$ with a single element
$p:U\rightarrow X$.
We obtain the functor
$\rho ^+\Uu: \Delta ^+ \rightarrow \Xx
^o$ defined by
$$
\rho ^+(\Uu )(p) := U \times _X \ldots \times _XU \;\;\; (p+1 \;\;
\mbox{times}).
$$
In particular
we have $\rho ^+ (\Uu )(\iota ) = X$. We have the left Quillen
presheaf $\rho ^+ (\Uu )^{\ast} \mm$ on $(\Delta ^+)^o$. We shall call
$\rho (\Uu )^{\ast} \mm$ its restriction to $\Delta$. By Corollary
\ref{pqcalclim} (this is where we strictify our descent data) we have
$$
\lim _{\leftarrow , \Delta} \rho (\Uu )^{\ast}L(\mm )
\cong L(Sect ^{\rm eq}(\Delta , \int _{\Delta}\rho (\Uu )^{\ast}\mm )).
$$
%%%%%%%% the following was changed in v3 %%%%%%%%%
To prove that $L(\mm )$ is a stack, it suffices, by Proposition
\ref{critere7} ((a) + (b)), to prove (b), that the restriction morphism
$$
L(\mm (X))\rightarrow \lim _{\leftarrow , \Delta} \rho (\Uu )^{\ast}L(\mm )
$$
is essentially surjective.
Indeed condition (a) of \ref{critere7}
is hypothesis (5) of the present statement.
%%%%%%%% end of change %%%%%%
So it suffices to prove that
$$
L(\mm (X))\rightarrow
L(Sect ^{\rm eq}(\Delta , \int _{\Delta}\rho (\Uu )^{\ast}\mm ))
$$
is an equivalence.
% the original prints $M(X)$ here, for $\mm (X)$
Since we have $M(X)=\rho (\Uu )^{\ast}\mm (\iota )$, the question concerns
only the restriction $\rho ^+(\Uu )^{\ast}\mm$ to $\Delta ^+$, and one can
apply Corollary \ref{dpq2} above.
One equips
$$
Sect (\Delta , \int _{\Delta}\rho (\Uu )^{\ast}\mm )
$$
with its cmc structure of type (II) guaranteed by hypothesis (0).

Hypotheses (0) and (1) imply
immediately conditions (o) and (i) of Corollary \ref{dpq2}.


For hypothesis (ii) of \ref{dpq2}, let
$$
\sigma \in
Sect ^{\rm eq}(\Delta , \int _{\Delta}\rho (\Uu )^{\ast}\mm )_f
$$
(this is our descent datum). We shall descend it to $\varphi _{\ast}
(\sigma )$.
We shall show that the adjunction morphism
$$
\varphi ^{\ast} \varphi _{\ast} (\sigma )\rightarrow \sigma
$$
is an equivalence, after which \ref{dpq2} will apply to give the
desired result.



Above $(\Delta ^+\times \Delta ^+)^o$ we have the left Quillen
presheaf $\nn$ defined by
$$
\nn (a,b) := \mm (\rho ^+ (\Uu )(a) \times _X \rho ^+ (\Uu )(b)).
$$
There is a functor
$$
\gamma : \Delta ^+ \times \Delta ^+ \rightarrow
\Delta ^+
$$
of concatenation of ordered sets ($\iota$ is the empty set),
and we have
$$
\nn = \gamma ^{\ast} (\rho (\Uu )^{\ast}\mm ).
$$
With the section $\sigma$ we obtain three sections of $\int _{\Delta \times
\Delta} \nn$, denoted
$$
p_1^{\ast}\sigma , \;\;\;
p_2^{\ast}\sigma , \;\;\;
\gamma ^{\ast}\sigma .
$$
The third one, $\gamma ^{\ast}\sigma$, is just the pull-back of $\sigma$
along
$\gamma$; and for example $p_1^{\ast}(\sigma )(a,b)$ is the image of
$\sigma (\rho (\Uu )(a))$ under the morphism
$$
\mm (\rho (\Uu )(a))\rightarrow \nn (a,b) = \mm (\rho (\Uu )(a)\times
_X\rho (\Uu
) (b))
$$
of restriction via the first projection of the product.
We have morphisms
$$
p_1^{\ast}\sigma \rightarrow \gamma ^{\ast} \sigma \leftarrow p_2^{\ast} \sigma
$$
and our hypothesis that $\sigma$ is an $eq$-section implies that these
morphisms are weak equivalences (object by object above
$\Delta
\times \Delta$).

We equip
$Sect (\Delta \times \Delta ,
\int _{\Delta \times \Delta}\nn )$
with a cmc structure which we shall call {\em of type (I-II)} (which means
{\em roughly} ``type (I) in the first variable and type (II) in the
second'') in the following way. For $p\in \Delta$ we equip
$$
Sect (\{ p\} \times \Delta   ,
\int _{\{ p\} \times \Delta   }\nn |_{\{ p\} \times \Delta })
$$
with its structure of type (II) generated by cofibrations (by
hypothesis (0)). These cmc's form a left Quillen presheaf on
$\Delta$ each of whose values is generated by cofibrations. The category of
sections of this Quillen presheaf---which we equip with its structure of type
(I) of Theorem \ref{sectionscmf}---is exactly $Sect (\Delta \times
\Delta , \int _{\Delta \times \Delta}\nn )$.

In this structure of type (I-II), the cofibrations are the morphisms which,
restricted to each $\Delta \times \{ p\}$, are cofibrations for the
structure of type (I) (i.e.\ of HBKQ type). On the other hand, the
fibrations are the morphisms which, restricted to each $\{ p\} \times
\Delta$, are fibrations for the structure of type (II) (i.e.\ of
Jardine--Brown--Heller type).

With these structures, the functor (``pull-back by the first projection'')
$$
p_1^{\ast} : Sect (\Delta , \int _{\Delta} \rho (\Uu )^{\ast} \mm )
\rightarrow
Sect (\Delta \times \Delta ,
\int _{\Delta \times \Delta}\nn )
$$
preserves cofibrations, i.e.\ it is a left Quillen functor.
Its right adjoint, which exists by (0), and which we shall denote
$$
p _{1,\ast}: Sect (\Delta \times \Delta ,
\int _{\Delta \times \Delta}\nn )\rightarrow
Sect (\Delta , \int _{\Delta} \rho (\Uu )^{\ast} \mm )
$$
is a right Quillen functor. It is
the analogue of $\varphi _{\ast}$ for the situation relative to the
first projection.

%%% corrected
%%% the following sentence was false in v1; the correct
%%% statement requires an additional hypothesis (4) in the theorem.
%%Condition (3) of the theorem implies that the natural
%%morphism
%%$$
%%\varphi ^{\ast} \varphi _{\ast} (\sigma )\stackrel{\cong}{\rightarrow}
%%p _{1,\ast}p_2^{\ast} (\sigma )
%%$$
%%is an equivalence.
%%%
%%%% here is the corrected version in v2:

We shall prove the following assertion:
the natural
morphism
$$
\varphi ^{\ast} \varphi _{\ast} (\sigma )\stackrel{\cong}{\rightarrow}
p _{1,\ast}p_2^{\ast} (\sigma )
$$
is an equivalence. For this, we shall use
conditions (3) and (4) of the theorem.
The sentence corresponding to this assertion in the first version
of the paper was incorrect; it is to correct the argument that
hypothesis (4) had to be added. For the proof, we begin by recalling that
$$
\varphi _{\ast} ( \sigma ) = \lim \, r_{\ast} (\sigma )
$$
where
$$
r_{\ast} (\sigma ) \in \mm (X) ^{\Delta} .
$$
The limit in question is homotopical (cf.\ the discussion at the beginning
of this chapter).


We have an arrow
$$
\ell : \varphi ^{\ast} \lim _{\Delta}\, r_{\ast} (\sigma ) \rightarrow
holim _{\Delta} \varphi ^{\ast} r_{\ast} (\sigma )
$$
and condition (4) of the theorem means that $\ell$
is an equivalence. Here,
$$
\varphi ^{\ast}r_{\ast} (\sigma ) \in
Sect (\Delta , \int _{\Delta}\rho (\Uu )^{\ast} \mm )^{\Delta},
$$
and the occurrence of $\Delta$ along which the $holim$ is taken is the one
which appears as an exponent
(and which corresponds to the variable denoted $b$ above).

On the other hand, condition (3) of the theorem means that the
natural arrow
$$
\varphi ^{\ast}r_{\ast} (\sigma ) \rightarrow
r_{1,\ast} p_2^{\ast} (\sigma )
$$
is an object-by-object equivalence above
$\Delta \times \Delta$, where
the notation $r_{1,\ast}$ denotes the morphism induced by
the adjoints of the restrictions
$$
r_{1,\ast} :
Sect (\Delta \times \Delta , \int _{\Delta \times \Delta}\nn )
\rightarrow
Sect (\Delta , \rho (\Uu )^{\ast} \mm )^{\Delta}.
$$
To see this, fix $(a,b)\in \Delta \times \Delta$. Then the arrow
in question at the point $(a,b)$ is the natural arrow between the two
composites in the following diagram:
$$
\begin{array}{ccccc}
\nn (a,b) = & \mm (\rho (\Uu )(a) \times _X \rho (\Uu )(b) )
& \stackrel{r_{1,\ast}}{\rightarrow} & \mm (\rho (\Uu )(a)) &
= \rho (\Uu )^{\ast} \mm (a) \\
& {\scriptstyle p_2^{\ast}}\uparrow  && \uparrow {\scriptstyle \varphi ^{\ast}}
& \\
\rho (\Uu )^{\ast}\mm (b) = & \mm (\rho (\Uu )(b)) &
\stackrel{r_{\ast}}{\rightarrow} & \mm (X) & = \rho (\Uu )^{\ast}(\iota ).
\end{array}
$$
Condition (3) of the theorem says exactly that this arrow is an
equivalence.


We show below (as in the initial version)
that $p _{1,\ast}p_2^{\ast} (\sigma )$ is the
homotopy limit (along $\Delta$) of
$r_{1,\ast} p_2^{\ast} (\sigma )$. Hence the equivalence given above
by condition (3) implies that
$$
p _{1,\ast}p_2^{\ast} (\sigma )\cong
holim \varphi ^{\ast} r_{\ast} (\sigma ).
$$
The arrow in question in the assertion is therefore equivalent to
the arrow $\ell$.

On the other hand, the $holim$ in
$Sect (\Delta , \rho (\Uu )^{\ast} \mm )$ are computed object by object
above $\Delta$. One can see this using the structure (I)
of Theorem \ref{sectionscmf}. Hypothesis (4) implies that
$\varphi ^{\ast}$, composed with any restriction to an object
of $\Delta$, commutes with homotopy limits. It follows that $\varphi ^{\ast}$
commutes
with homotopy limits.
We conclude that $\ell$ is an equivalence, which gives the assertion.


For the rest of the present proof we take up
the first version of the paper. Further on, in
\S 20 and \S 21, we shall have to verify hypothesis
(4) of the theorem when using it.
%%% end of the correction in v2



%%% end of correction



Let ${\bf f}$ be the ``fibrant replacement'' functor for the structure
of type (I-II) on the sections of $\nn$. The functor
$p _{1,\ast} \circ {\bf f}$ is invariant under equivalence, so we
obtain equivalences
$$
p _{1,\ast} {\bf f}(p_1^{\ast} \sigma )\rightarrow
p _{1,\ast} {\bf f}(\gamma ^{\ast}\sigma ) \leftarrow
p _{1,\ast} {\bf f}(p_2^{\ast} \sigma ).
$$
Under the hypothesis that $\sigma$ is fibrant (for the structure of
type (II)), we
have that $p_2^{\ast}\sigma$ is fibrant for the structure of type (I-II),
and hence
$$
p _{1,\ast}p_2^{\ast} (\sigma )
\rightarrow
p _{1,\ast} {\bf f}p_2^{\ast} (\sigma )
$$
is an equivalence.
Note that this implies that $p _{1,\ast}p_2^{\ast} (\sigma )$ is a
homotopy limit of $r_{1,\ast}p_2^{\ast} (\sigma )$.
On the other hand we have the adjunction morphism
$$
\varphi ^{\ast} \varphi _{\ast} (\sigma )
\rightarrow \sigma
$$
and also
$$
\sigma \rightarrow p _{1,\ast} {\bf f} p_1^{\ast} \sigma .
$$
Their composite is the morphism
on the left in the square $(\ast )$
$$
\begin{array}{ccc}
\varphi ^{\ast} \varphi _{\ast} (\sigma )&\rightarrow &
p _{1,\ast} p_2^{\ast} (\sigma )\\
\downarrow && \downarrow \\
p _{1,\ast} p_1^{\ast} (\sigma )&\rightarrow &
p _{1,\ast} \gamma ^{\ast}(\sigma ).
\end{array}
$$
The other morphisms are deduced from the preceding morphisms.
It is a not altogether obvious fact that this square commutes.
To prove it, set $\eta :=  \varphi _{\ast} (\sigma )$
and consider
first the square $(\ast \ast )$
$$
\begin{array}{ccc}
t^{\ast} \eta & \rightarrow & p _2^{\ast} \sigma \\
\downarrow && \downarrow \\
p_1^{\ast} \sigma & \rightarrow & \gamma ^{\ast} \sigma ,
\end{array}
$$
where
$$
t^{\ast} \eta := p_1^{\ast} \varphi ^{\ast} \eta = p_2^{\ast} \varphi
^{\ast} \eta .
$$
This square $(\ast \ast )$ commutes. To see this, let us denote by $P_a$ the
fibre product of $a+1$ copies of $U$ over $X$ corresponding to
$a\in \Delta$; and let us denote by $[x\rightarrow y]^{\ast} $ and $[x\rightarrow
y]_{\ast}$ the restriction functors and their adjoints. Note that we have
$$
P_{\gamma (a,b)} = P_a\times _XP_b
$$
and
$$
\eta = \lim _{\leftarrow , a} [P_a\rightarrow X]_{\ast}(\sigma (a)).
$$
Let us denote by $
f_a: \eta \rightarrow [P_a\rightarrow X]_{\ast}(\sigma (a)).
$ the corresponding structural morphisms.
The transition morphisms in the limit come for example
from the structural morphisms of $\sigma$ such as
$$
[P_a\times _XP_b\rightarrow P_a]^{\ast} (\sigma (a)) \rightarrow
\sigma (\gamma (a,b)),
$$
which gives
$$
[P_a\rightarrow X]_{\ast} (\sigma (a))\rightarrow
[P_a\times P_b\rightarrow X]_{\ast}
[P_a\times _XP_b\rightarrow P_a]^{\ast} (\sigma (a))
$$
$$
\rightarrow
[P_a\times P_b\rightarrow X]_{\ast}\sigma (\gamma (a,b)).
$$
We obtain that the composite
$$
\eta \stackrel{f_a}{\rightarrow}
[P_a\rightarrow X]_{\ast} (\sigma (a))\rightarrow
\rightarrow
[P_a\times P_b\rightarrow X]_{\ast}\sigma (\gamma (a,b))
$$
is equal to the morphism $f_{\gamma (a,b)}$.
The composite in question is the adjoint of
the component at $(a,b)$ of the lower composite in the square $(\ast \ast )$.
It follows that this composite
is equal to the natural diagonal morphism
$$
t^{\ast} \eta\rightarrow \gamma ^{\ast} \sigma
$$
whose component at $(a,b)$ is the adjoint of $f_{\gamma (a,b)}$.
By symmetry, the square $(\ast \ast )$ commutes.

Applying to the square $(\ast \ast )$ the operation
$p_{1,\ast }$ we obtain a square which one can compose with the square
of naturality of the adjunction morphisms
$$
\begin{array}{ccc}
\varphi ^{\ast} \eta & \rightarrow & p_{1,\ast } p_1^{\ast} \varphi ^{\ast}
\eta \\
\downarrow && \downarrow \\
\sigma &\rightarrow & p_{1,\ast } p_1^{\ast} \sigma
\end{array}
$$
to show that the square $(\ast )$ commutes.

Applying the operation ``fibrant replacement'' in the middle of the bottom
row we obtain (by naturality of the natural transformation $u\rightarrow
{\bf f}u$) the square $(\ast )'$
$$
\begin{array}{ccc}
\varphi ^{\ast} \eta &\rightarrow &
p _{1,\ast} p_2^{\ast} (\sigma )\\
\downarrow && \downarrow \\
p _{1,\ast} {\bf f} p_1^{\ast} (\sigma )&\rightarrow &
p _{1,\ast} {\bf f} \gamma ^{\ast}(\sigma )
\end{array}
$$
which commutes. Now here we already know that the morphisms at the top, on
the right
and at the bottom are equivalences. By ``three for the price of two'', the
morphism on the left is an equivalence. Recall that it is the
composite
$$
\varphi ^{\ast} \varphi _{\ast} (\sigma ) =\varphi ^{\ast} \eta
\rightarrow \sigma \rightarrow p _{1,\ast} {\bf f} p_1^{\ast} (\sigma ).
$$
In sum, to prove that our adjunction morphism
$$
\varphi ^{\ast} \varphi _{\ast} (\sigma ) \rightarrow \sigma
$$
is an equivalence, it suffices to
prove that the morphism
$$
\sigma \rightarrow p _{1,\ast} {\bf f} p_1^{\ast} (\sigma )
$$
is one.


This last fact belongs to more or less standard homological or homotopical
algebra. We indicate the argument briefly, referring to Illusie
\cite{Illusie}
\footnote{
The reader should beware that the meaning of
our exponent $+$ (e.g.\ $\Delta
^+$) differs from that of \cite{Illusie}---it is
rather exactly the opposite.}
(who himself refers to ``secret
papers'' of Deligne---but the reader will be able to prove the result in
question on his own). Let us set $$
P_m:= \rho (\Uu )(m) = \Uu \times _X\ldots \times _X \Uu .
$$
We shall use the notations $[x\rightarrow y]^{\ast}$ and
$[x\rightarrow y]_{\ast}$ for the restriction functor and its adjoint
between $\mm (y)$ and $\mm (x)$,
corresponding to a morphism $x\rightarrow y$. On the other
hand we work in $\Xx /X$, so we may suppress the subscript $X$ for the
fibre product over $X$.
With these notations, for $(a,b)\in \Delta \times \Delta$ we have
$$
(p_1^{\ast} \sigma )(a,b)=
[P_a \times P_b \rightarrow P_a]^{\ast} \sigma (a).
$$
We obtain the cosimplicial object of $\mm (P_a)$, with (co)simplicial index
$b$
$$
C_{a,b}:= [P_a \times P_b \rightarrow P_a]_{\ast}
[P_a \times P_b \rightarrow P_a]^{\ast}\sigma (a).
$$
We have
$$
(p _{1,\ast} {\bf f}p_1^{\ast} \sigma )(a)\cong
holim _{b\in \Delta} C_{a,b}.
$$
Let us denote by $\epsilon : P_a\times P_0 \rightarrow P_a$
the projection. One can decompose
the projection $P_a\times P_b\rightarrow P_a$ as a product of $b+1$ copies of
$\epsilon$. We denote by $\epsilon ^{\ast}$ and $\epsilon _{\ast}$ the
functorialities for $\mm$ with respect to
$\epsilon$. Using the formula
given by hypothesis (2) of the present theorem, we have
$$
C_{a,b} = (\epsilon _{\ast} \epsilon ^{\ast} )^{b+1}\sigma (a).
$$
In particular, this cosimplicial object of $\mm (P_a)$ is the {\em standard
cosimplicial resolution of $\sigma (a)$ provided by the pair of adjoint
functors $\epsilon _{\ast}, \epsilon ^{\ast}$}---cf.\ Illusie \cite{Illusie}
\S 1.5. Theorem 1.5.3 of {\em loc cit.} implies that
the augmented cosimplicial object
$$
\epsilon ^{\ast} \sigma (a) \rightarrow \epsilon ^{\ast} C_{a,-}
$$
of $\mm (P_a\times P_0)$ is ``homotopically trivial''.
By \cite{Illusie} (1.1.5 and the description for 1.5.3) this means
that there is a morphism
$$
\epsilon ^{\ast} C_{a,-}\rightarrow \epsilon ^{\ast} \sigma (a)
$$
such that the composite
% the original prints $\epsilon^{\ast}C_{a,-}$ as the source here
$$
\epsilon ^{\ast} C_{a,-}\rightarrow
\epsilon ^{\ast} C_{a,-}
$$
is homotopic (in the sense of \cite{Illusie} 1.1.5) to the identity
(the other composite being itself the identity).

There is a section
$$
s:P_a\rightarrow P_a\times P_0
$$
(diagonal in the last variable), and applying $s^{\ast}$ to the augmented
cosimplicial object above, we obtain (with the formula
$s^{\ast} \epsilon ^{\ast} = 1$) that
the augmented cosimplicial object
$$
\sigma (a) \rightarrow C_{a,-}
$$
of $\mm (P_a)$ is homotopically trivial.
It follows (since $holim$ turns the homotopies of \cite{Illusie} 1.1.5
into Quillen homotopies \cite{Quillen} for the cmc structure of $\mm
(P_a)$) that the morphism
$$
holim _{b\in \Delta}\sigma (a)\rightarrow  holim _{b\in \Delta}
C_{a,b}
$$
admits an inverse up to (Quillen) homotopy. This implies
that it is a weak equivalence. Note that
$$
\sigma (a)\rightarrow holim _{b\in \Delta}\sigma (a)
$$
is an equivalence. For this we refer to the result of Hirschhorn
(\cite{Hirschhorn}
Theorem 19.5.1) which implies that the $holim$ of a constant diagram (in
any cmc satisfying (0), for example), taken over a category
whose nerve is contractible (as is the case for $\Delta$), is equivalent
to the value taken on the diagram.

We conclude that we have an equivalence
$$
\sigma (a)\stackrel{\cong}{\rightarrow} holim _{b\in \Delta}
C_{a,b}
$$
which completes the proof.
\eop

{\em Remark:}
One might think that if $M$ is a fixed closed model category
admitting arbitrary limits and colimits, then the constant presheaf
$\underline{M}$ with values in $M$ satisfies the hypotheses of the
theorem for any topology $\Gg$. Yet $L(\underline{M})$, which is
the constant Segal $1$-prestack with values $L(M)$, is not in general
a stack. The reason is that $L(\mm )$ is not compatible with disjoint
sums.

{\bf Exercise:} Rewrite the statement and the proof of the theorem
\ref{dansleschamps} with Reedy type structures instead of type (II);
without the hypothesis that all objects of the $\mm (y)$ are cofibrant; and
with the morphisms
of hypotheses (2) and (3) assumed to be
weak equivalences only
when the objects concerned are cofibrant and fibrant.

We point out that the authors have not done this exercise.

\subnumero{Comparison with SGA 4}
We compare here our result of the present chapter with that of
Exposé Vbis of SGA 4 \cite{SGA4b} by Saint-Donat (after ``succinct
notes'' of Deligne).


The first remark to make is that we consider the
``unstable'' situation via our
arbitrary cmc's, whereas \cite{SGA4b} is devoted to the descent of
complexes, which
corresponds to stable homotopy. Nevertheless, the framework is very close,
and the application which
motivated us is precisely the example of complexes. To make this
transcription, note that
in \cite{SGA4b} what is involved is a family of ringed toposes indexed
by the base
category; one can recover the ``family of cmc's'' framework by taking the
family of
cmc's of complexes of
modules over the structure rings in the fibre toposes (we do not enter
into the details of a
potential proof of the present assertion). The object
$D^+(\underline{\Gamma}(E), A)$
of \cite{SGA4b} Vbis 2.2.7 is the analogue of our $Sect (\Delta , \int
_{\Delta}\mm )$; and
the full subcategory in question in
\cite{SGA4b} Vbis 2.2.7 (of sections whose cohomologies are
cocartesian) corresponds to our $Sect ^{\rm eq}(\Delta , \int
_{\Delta}\mm )$.

In the framework of the descent of complexes (i.e.\ for the application to
\S 21) one can very well
substitute the argument of \cite{SGA4b} for our argument of the present
chapter. However, it is
to be noted that our Theorem \ref{dansleschamps} also comprises a part
``strictification of
descent data'', made concrete in the use of Corollary
\ref{pqcalclim}.
This part has no counterpart in \cite{SGA4b}, where one starts
from a ``descent
datum'' which is already a complex in the topos of sections,
that is, a strict family
of complexes. The language of
$\infty$-categories, which allows us to introduce
the notion of weak descent datum, was not available at the time.

As for the method
for descending a strict datum, the argument of
\cite{SGA4b} for the case
of complexes is a dévissage by truncation of the complex, allowing one to
treat the cohomology objects one at a time. Starting from a cocartesian
section $\sigma$, one applies the
``direct image'' functor which we call $\varphi _{\ast}(\sigma )$
(in the notation of
\cite{SGA4b} the descent datum is $F^{\cdot}$ and its direct image is
${\bf R}^+ (\overline{\theta}_{\ast})(F^{\cdot})$). It is a matter of showing
that $\varphi
^{\ast}\varphi_{\ast} \sigma  $ is equivalent to $\sigma$; which is
done on one page for the
case of complexes in \cite{SGA4b} Vbis Prop. 2.2.7.

This method would also be conceivable for the ``unstable'' case,
provided one had
a sufficient theory of the ``Postnikov tower''.

The major part of \cite{SGA4b} is devoted to the search for conditions
guaranteeing that
the inverse image ${\bf L}^+(\overline{\theta}^{\ast})$ (which we call
$L(\varphi ^{\ast})$) is
fully faithful. Our Lemma \ref{dpq1} is the analogue of the remark
(\cite{SGA4b} Vbis Prop. 2.2.7).
Then (for our Corollary \ref{dpq2},
whose proof is to be found
after \ref{dpq1})
we used condition \ref{dpq2} (ii), which
corresponds to the effectivity
of descent data.
This makes the argument of our \ref{dpq2} easier; and one can do it
because in any case
we need a specific argument to obtain the effectivity
of
\ref{dpq2} (ii). The statement of \ref{dpq2} is to be compared with the
statement of
\cite{SGA4b} Vbis Prop. 2.2.7.

Let us look finally at the various hypotheses in our Theorem
\ref{dansleschamps}.
The hypothesis on the existence of sufficiently many disjoint sums is to
be compared with
\cite{SGA4b} Vbis 3.0.0. One can compare condition (3) with
\cite{SGA4b} Vbis 3.2.
Here one may remark that Saint-Donat uses an arbitrary family of
base changes in order to
``test'', whereas our strategy is to test with the
base change which one
uses for the descent.

The bisimplicial technique of the proof of \ref{dansleschamps}
is akin to the
bisimplicial technique for full faithfulness of SGA 4 (\cite{SGA4b}
Vbis \S 2.3, see also
3.3.1(a)).

In sum, the overall strategy, which consists in
descending $\sigma$ by taking
$\varphi _{\ast}(\sigma )$,
is common, but the problem of proving that this really answers the
question is for us more
complicated, since we do not have a ``Postnikov-style'' recurrence
argument at our disposal. This then requires
a specific argument (proof of \ref{dansleschamps}); and with
this one can bypass the
problem of the full faithfulness of $L(\varphi ^{\ast})$ (\ref{dpq2}).
Naturally, the
arguments of \cite{SGA4b} for full faithfulness are found again in
our specific argument
for the effectivity of
\ref{dansleschamps}.
